\section*{Chapitre \ref{ch:ml:data-and-feature-stores} --- Données et feature stores}

\begin{solution}{exo:ml:data-and-feature-stores:1}
Le \omterm{def:ml:data-and-feature-stores:store}{magasin en ligne}~: transactions horodatées dans $(09{:}59{:}37.300, 10{:}00{:}07.300]$. La chaîne des barres~: transactions comprises dans les trente barres
d'une seconde se terminant à 10:00:08, soit $(09{:}59{:}38, 10{:}00{:}08]$~: elle manque 0.7~s du passé et inclut 0.7~s du futur.
\end{solution}

\begin{solution}{exo:ml:data-and-feature-stores:2}
$1 - 0.999^n\ge0.95$ donne $n\ge\ln0.05/\ln0.999 = 2\,994$ décisions~: environ 3\,000, bien plus qu'un contrôle sur cinq décisions, ce qui explique
pourquoi les contrôles de parité doivent aussi tourner sur des journées entières, en batch.
\end{solution}

\begin{solution}{exo:ml:data-and-feature-stores:3}
Les variables de volume et de comptage sont des sommes d'entiers, représentées exactement en float32 jusqu'à $2^{24}$~; les déséquilibres et le
mid pondéré sont des rapports à 24 bits de précision en float32, contre 53 en float64~: ils changent donc au septième ou au huitième
chiffre significatif.
\end{solution}

\begin{solution}{exo:ml:data-and-feature-stores:4}
La fin d'une barre tombe jusqu'à une seconde après la décision~: pour une prévision à une seconde, l'intervalle fuité représente une grande part
de la fenêtre cible, et le modèle de recherche a appris à l'exploiter~; pour une prévision à cinq secondes, elle représente au plus un cinquième de la fenêtre, et
l'information est diluée.
\end{solution}

\begin{solution}{exo:ml:data-and-feature-stores:5}
Le message qui suit une transaction est en général un ordre à cours limité ajouté ou annulé au fond du carnet d'ordres~: son effet sur la cible est faible et
il ajoute du bruit aux variables sur lesquelles s'appuie le modèle, de sorte que le modèle de recherche est légèrement moins bon qu'un modèle entraîné sur des variables propres. Si
le message suivant était typiquement celui qui déplace le mid (carnet mince, marché rapide), le biais d'anticipation gonflerait plutôt
le backtest~; le signe de l'effet d'une fuite n'est pas connaissable à l'avance.
\end{solution}

\begin{solution}{exo:ml:data-and-feature-stores:6}
Des moments égaux ne signifient pas des valeurs égales~: des barres d'une seconde décalent les valeurs dans le temps sans modifier leur distribution, et un skew
qui touche certaines décisions peut laisser intactes les moyennes journalières. Il faut comparer les valeurs elles-mêmes, décision par décision.
\end{solution}

\begin{solution}{exo:ml:data-and-feature-stores:7}
Entraîné sur le rejeu du moteur de production, l'IC en production du modèle vaut 0.411, le résultat propre, contre 0.360 pour le modèle entraîné sur
des variables de barres~: tout ce que coûtait le skew est récupéré. Le prix à payer est de rejouer tout l'historique des messages dans le moteur en ligne
pour construire l'\omterm{def:ml:why-financial-machine-learning-is-different:sets}{ensemble d'entraînement}, ce qui est plus lent que la chaîne des barres, et c'est précisément ce que la définition partagée d'un \omterm{def:ml:data-and-feature-stores:store}{feature store} rend routinier.
\end{solution}

\begin{solution}{exo:ml:data-and-feature-stores:8}
L'erreur d'une somme en float64 de $n$ termes, dans n'importe quel ordre, est majorée environ par $n\,u\sum|x_i|$ avec $u = 2^{-53}$~; utiliser une tolérance
proportionnelle à $\sum|x_i|$ multiplié par un petit multiple de $n\,u$, ou rendre les deux calculs exacts (entrées entières, sommation
compensée, livre~4, chapitre~25) et exiger l'égalité.
\end{solution}

\begin{solution}{pb:ml:data-and-feature-stores:1}
\textbf{Partie I.}
\begin{enumerate}
\item Nom, entrée, agrégation, fenêtre, horloge et version~; sous forme de données, pour que les deux magasins, la documentation et les tests lisent la
même déclaration.
\item Hors ligne~: vectorisé sur un historique à partir de sommes cumulées et de jointures as-of~; en ligne~: fenêtres glissantes mises à jour à chaque
message. Elles coïncident exactement pour les 7\,714 décisions.
\item L'instant du dernier message lu~; pour une donnée fondamentale, son instant de livraison à la société.
\item L'âge de l'entrée la plus récente d'une valeur~; les transactions publiées avec 200~ms de retard ont rendu périmées les variables de transaction.
\end{enumerate}
\textbf{Partie II.}
\begin{enumerate}[resume]
\item Une avance d'un événement (toutes les variables de carnet, un peu les fenêtres), barres d'une seconde (toutes les variables, surtout les dernières valeurs du
carnet), stockage en float32 (les variables de ratio, sans effet visible), transactions publiées en retard (les trois variables de transaction).
\item Les deux premiers sont des erreurs de recherche~; les deux derniers sont des conditions de production.
\item 100~\%, 96~\%, 92~\% et 100~\% des décisions (à la tolérance $10^{-9}$).
\item À $10^{-9}$, il signale le stockage en float32, pourtant inoffensif~; à $10^{-6}$, il l'ignore et signale tout de même les autres.
\end{enumerate}
\textbf{Partie III.}
\begin{enumerate}[resume]
\item Backtest et en production~: propre 0.411 et 0.411~; un événement d'avance 0.395 et 0.418~; barres d'une seconde 0.438 et 0.360~; float32 0.411
et 0.411~; transactions tardives 0.411 et 0.411.
\item Ses variables ont vu une partie de la fenêtre cible~; le modèle les a pondérées~; en production, elles ne portent aucune information de ce type.
\item Les mêmes skews modifient l'IC de moins de 0.01.
\item Un IC en production nettement inférieur à celui du backtest dès le premier jour, stable ensuite~: la signature d'un écart de chaîne,
et non d'une dégradation.
\end{enumerate}
\textbf{Partie IV.}
\begin{enumerate}[resume]
\item \emph{Mêmes variables, deux fois.} La recherche sur barres d'une seconde gonfle l'IC du backtest de la prévision à une seconde jusqu'à 0.438 et
donne 0.360 en production contre 0.411 pour la chaîne propre, le P\&L par décision tombant de 0.089 à 0.078 ticks~; un
contrôle de parité sur cinq décisions détecte à chaque fois chacun des quatre skews introduits.
\item Le décalage des barres~: il gonfle la recherche et dégrade la production, et rien, pris isolément dans l'une ou l'autre, ne le révèle.
\item Une nouvelle version calculée à côté de l'ancienne, rétro-remplie hors ligne, comparée en ligne, et basculée seulement lorsque les modèles sont
réentraînés sur elle.
\item Arrêter de trader avec les modèles touchés ou revenir à une version antérieure, identifier la variable et sa cause, corriger, recalculer, et réentraîner si les
données d'entraînement étaient fausses.
\item Par la \omterm{def:ml:alternative-data-pipelines:ingest}{chaîne d'ingestion} du chapitre~15, avec les instants de première observation comme temps de connaissance, puis comme définitions ordinaires.
\item Définitions, versions, historique des contrôles de parité et journaux de \omterm{def:ml:data-and-feature-stores:fresh}{fraîcheur des variables}, pour chaque variable que le modèle utilise.
\item Du temps d'ingénierie et de la discipline~; cela vaut la peine dès que deux modèles partagent des variables ou qu'un modèle passe des ordres en production.
\item Que le modèle voit en production ce qu'il a vu à l'entraînement, à l'instant où il décide.
\end{enumerate}
\end{solution}

\begin{solution}{iq:ml:data-and-feature-stores:1}
Un écart entre les variables sur lesquelles un modèle a été entraîné et celles qu'il reçoit en production~: chemins de code différents, horloges ou barres différentes,
entrées tardives ou manquantes, arithmétique différente, biais d'anticipation dans la jointure de recherche.
\iqlookfor{la définition et trois causes distinctes.}
\end{solution}

\begin{solution}{iq:ml:data-and-feature-stores:2}
Définitions déclaratives~; un moteur en ligne qui met à jour des fenêtres glissantes à chaque message avec une latence bornée~; un \omterm{def:ml:data-and-feature-stores:store}{magasin hors ligne}
qui rejoue le même moteur ou un calcul vectorisé équivalent sur l'historique, horodaté par les temps de connaissance~; une récupération
point-in-time pour l'entraînement~; suivi quotidien de la parité et de la fraîcheur~; versionnement.
\iqlookfor{définitions partagées, hors ligne point-in-time, parité et fraîcheur.}
\end{solution}

\begin{solution}{iq:ml:data-and-feature-stores:3}
Rejouer les messages de la période en production dans la chaîne de recherche et comparer les variables à ce que la production a journalisé~; réentraîner sur
les variables rejouées par la production et voir si le backtest tombe à l'IC en production~; contrôler les journaux de \omterm{def:ml:data-and-feature-stores:fresh}{fraîcheur des variables}.
\iqlookfor{une comparaison par rejeu avant toute explication de modélisation.}
\end{solution}

\begin{solution}{iq:ml:data-and-feature-stores:4}
Pour chaque couple (entité, instant de décision), prendre la dernière valeur de chaque variable dont le temps de connaissance est antérieur ou égal à l'instant de décision~; mis en œuvre
par des jointures as-of sur des temps triés, jamais par une jointure sur des clés calendaires.
\iqlookfor{temps de connaissance et règle as-of.}
\end{solution}

\begin{solution}{iq:ml:data-and-feature-stores:5}
Enregistrer la dernière arrivée de chaque entrée et l'entrée la plus récente de chaque variable~; déclencher une alerte lorsque leur âge dépasse une borne
fixée par la fenêtre de la variable et l'horizon du modèle~; faire refuser au modèle, ou le faire se replier, les variables périmées plutôt que de passer des ordres sur elles.
\iqlookfor{âge mesuré, bornes et comportement défini en cas de péremption.}
\end{solution}

\begin{solution}{iq:ml:data-and-feature-stores:6}
Exprimer les entrées en entiers (ticks, actions) ou en virgule fixe, de sorte que les sommes soient exactes~; ou utiliser le même ordre de sommation et une sommation
compensée dans les deux cas~; et tester sur un rejeu que chaque valeur coïncide exactement.
\iqlookfor{arithmétique exacte ou ordre contrôlé, et test par rejeu.}
\end{solution}
